
junho/2026

- verificar banca defesa

- programar defesa MSc
* junho/julho
* agosto/setembro
* outubro
* 30/10-01/10? - solicitacao aprovacao banca
* 05/10-16/10 - deposito
* 14/11/26 - defesa



16/04



% 12/03

% Reproduzir https://arxiv.org/pdf/2504.20872 e adicionar Interactive e Continual learning para comparar com os artigos

% https://arxiv.org/pdf/2504.11112



% Reproduzir SAMNet, UNet para entender os modelos de segmentação

% Preparar os slides para o Eldorado

% Preparar apresentação (10 minutos)


% ------------------------

% Usar FLIM, samNet, UNet --> Considerar ResNet
% Datasets -> BRaTS e Parasite

% Estudar e avaliar melhores estratégias de AL 
% Criar Slide para mostrar como os trabalhos atuais tem feito e escolhido as estratégias --- AL seleção de regiões

% 18/06
    ->  Mostrar visualmente antes vs depois (Apresentar imagens para cada decoder - e depois de aplicar AL também).
    -> Entender arquitetura de escolha de pesos (por que 3 a 5 imagens?)
    -> Melhorar como explorar o Active Learning (Não só selecionar o conjunto de imagens, como explorar melhor o AL (como explorar conjunto de pixels ao invés de imagens).
    -> Verificar a intervenção do especialista com marcador em imagens de alta confusão (entropia).
    -> Levantar datasets interessantes (mais informativos, menos informativos, com saliencia, sem saliencia, diversos) para aplicar o modelo.
    -> Levantar estado da arte para ver como os trabalhos estão aplicando AL (Entropia, Diversidade, etc..) - Conjunto de pixels mais relevantes? Aplicar também. (Levantar e apresentar como os trabalhos tem feito o uso de AL - POr exemplo, eles pegam as regiões mais relevantes? pixels com maiores incertezas?).
    
    